\paragraph{Network operators and their degree.}
We spell out the simultaneous argument, including the product over
all addresses.  Write $\partial_e$ for the finite difference in a
terminal pair bit $e$.  Products of these operators commute and
$\partial_e^2=0$.  The contraction rule identifies a forest product
with its contracted graph.  A cycle product is zero: after a
spanning path is contracted its last edge is internal and changes
neither a type nor any external incidence.  No contraction removes
the last vertex.

A network of primary auxiliary vertices acts on any of the symbols
by an operator in the $\partial_e$.  To see this directly, expand in
the bits incident to the auxiliary vertices.  Contract the chosen
forest edges.  A component with no terminal becomes an isolated
primary vertex, and its contribution vanishes by primary isolation.
A component meeting several terminals gives their forest difference.
If $Q$ is the grounded auxiliary Laplacian and $c_a$ is its attachment
column at terminal $a$, the degree-zero and degree-one terms are
\[
 \det(Q)\,1+
 \sum_{a<b}c_a^{\mathsf t}\operatorname{adj}(Q)c_b\,
                    \partial_{ab}.
\]
All other terms have difference degree at least two.  The displayed
coefficients are the rooted-forest minors, by Cauchy--Binet on the
incidence matrix.  They are polynomial identities, so apply to
singular $Q$ as well.  Distinct networks have no edges between their
auxiliary vertices.  Their operators multiply: expanding both
forests gives their union, and unions containing a terminal cycle
give zero by the preceding observation.  This also shows that the
operator is the same for both projected symbols.

Fix a target $x_*\ne0$.  Choose a basis
$\lambda_1,\ldots,\lambda_b$ of the dual address space with
$\lambda_i(x_*)=1$ for all $i$.  Such a basis is obtained from one
functional taking value one by adding it to a basis of the space
of functionals vanishing at $x_*$.  Group $i$ has three primary vertices, each activated
by $\lambda_i$.  In its physical three-dimensional coordinate space
let $e_i=(1,1,1)$ and choose $u_i,v_i$ spanning its orthogonal
complement.  The vector $e_i$ has norm one.  Prescribe the symmetric
matrix $Q$ by its bilinear matrix $I_b$ on the $e_i$ and
\[
 H(T,B)=\begin{pmatrix}0&T\\T^{\mathsf t}&B\end{pmatrix}
\]
on the remaining coordinates.  In physical coordinates each
$h_i=Qe_i=e_i$ is supported on its own group.  Attach this column to
a terminal activated by $\lambda_i$, which exists by hypothesis.
For an active set $J$ the sum of these active attachment columns is
$Q_J\mathbf1_J$.  Thus the prescribed principal matrix $Q_J$ really
is the grounded graph Laplacian on that address, including its
diagonal entries.

For a selected terminal pair $a,b$ active at $x_*$, let their
activations be $\alpha,\beta$.  Attach an extra column $c$, dual to
$u_1$, to both.  If $\alpha\ne\beta$, attach the same column to a
terminal of activation $\alpha+\beta$; it is absent at $x_*$.  The
extra columns sum to zero at every address since
$\alpha+\beta+(\alpha+\beta)=0$.  If $\alpha=\beta$, the two
columns already cancel.  Hence they do not alter $Q_J$.

For the singular network set $T=I+P$, where $P$ is the cyclic shift,
and give $B$ one diagonal entry one.  On every proper group subset,
the restricted cyclic shift is a union of paths, so $T_J$ is
invertible and $\det Q_J=\det(T_J)^2=1$.  On all groups, the kernels
of $T$ and $T^{\mathsf t}$ are generated by the all-one vector.
The one nonzero diagonal entry of $B$ forces the lower component of
a vector in $\ker H$ to vanish.  Thus $\ker Q$ is precisely the line
generated by $r=\sum_i u_i$, and
$\operatorname{adj}(Q)=rr^{\mathsf t}$.  The anchor columns pair
to zero with $r$, while $c^{\mathsf t}r=1$.  Only $a,b$ receive
active extra columns at $x_*$.  Consequently its network operator
$S_{x_*,e}$, for $e=ab$, has the expansions
\[
 S_{x_*,e}|_{x_*}=\partial_e+\text{terms of degree at least two},
 \qquad
 S_{x_*,e}|_x=1+\text{terms of positive degree}\quad(x\ne x_*).
\]
For $b=1$, $T=0$ and $B=(1)$ give exactly the same calculation.

For a varying network set $T=I+L$, with $L$ the consecutive lower
shift, and $B_\varepsilon=\varepsilon E_{bb}$ for
$\varepsilon\in\{0,1\}$.  Every $T_J$ is invertible and both
operators $V^\varepsilon_{x_*,e}|_x$ have constant coefficient one.
The upper-left block of $H^{-1}$ is
$T^{-\mathsf t}B_\varepsilon T^{-1}$.  Its change on the selected
columns is $((T_J^{-1})_{b1})^2$, interpreted as zero when an endpoint
is absent.  This is one exactly when the entire path is present,
which means $x=x_*$.  Anchor columns have zero change.  Therefore
\[
 \begin{aligned}
 (V^0_{x_*,e}+V^1_{x_*,e})|_{x_*}
       &=\partial_e+\text{terms of degree at least two},\\
 (V^0_{x_*,e}+V^1_{x_*,e})|_x
       &=\text{terms of degree at least two}\quad(x\ne x_*).
 \end{aligned}
\]

